\documentclass[11pt]{article}

\usepackage{amsmath}
\usepackage{amssymb}
\usepackage{geometry}
\usepackage{natbib}
\geometry{margin=1in}

\title{Fourier Analysis of Galactic Bars}
\author{}
\date{}

\begin{document}

\maketitle

\section*{1. Fourier Definition}

For a galaxy image in polar coordinates $(r, \theta)$:

\begin{equation}
A_m(r) = \frac{1}{I_0(r)} \int_0^{2\pi} I(r,\theta)\, e^{-im\theta} \, d\theta
\end{equation}

\subsection*{Discrete Form}

\begin{equation}
A_m(r) = \frac{\sum I_i(r)\cos(m\theta_i)}{\sum I_i(r)}
\end{equation}

\begin{equation}
B_m(r) = \frac{\sum I_i(r)\sin(m\theta_i)}{\sum I_i(r)}
\end{equation}

\subsection*{Amplitude}

\begin{equation}
|A_m(r)| = \sqrt{A_m^2 + B_m^2}
\end{equation}

\subsection*{Bar Mode ($m = 2$)}

\begin{equation}
A_2(r) = \frac{\sum I_i \cos(2\theta_i)}{\sum I_i}
\end{equation}

\begin{equation}
B_2(r) = \frac{\sum I_i \sin(2\theta_i)}{\sum I_i}
\end{equation}

\begin{equation}
\frac{A_2}{A_0} = \sqrt{A_2^2 + B_2^2}
\end{equation}

This quantity represents the \textbf{bar strength}.

\section*{2. Improvements for Real Data}

\subsection*{Deprojection}

Correct for inclination and rotate the galaxy to face-on orientation.

\subsection*{Elliptical Annuli}

\begin{equation}
r^2 = x'^2 + \frac{y'^2}{q^2}
\end{equation}

where $q = b/a$.

\subsection*{Phase}

\begin{equation}
\phi_2(r) = \frac{1}{2} \tan^{-1}\left(\frac{B_2}{A_2}\right)
\end{equation}

Bar region corresponds to approximately constant phase.

\section*{3. Bar Size Estimation}

\begin{equation}
R_{\text{bar}} = R(A_2^{\max})
\end{equation}

\section*{4. Error Estimation}

\begin{equation}
\sigma_{R_{\text{bar}}} = \sqrt{
\sigma_{\text{method}}^2 +
\sigma_{\text{MC}}^2 +
\sigma_{\text{bin}}^2
}
\end{equation}

\section*{5. Dynamical Ratio}

\begin{equation}
\mathcal{R} = \frac{R_{\text{corot}}}{R_{\text{bar}}}
\end{equation}

\begin{equation}
\sigma_{\mathcal{R}} =
\mathcal{R} \sqrt{
\left(\frac{\sigma_{R_{\text{corot}}}}{R_{\text{corot}}}\right)^2 +
\left(\frac{\sigma_{R_{\text{bar}}}}{R_{\text{bar}}}\right)^2
}
\end{equation}

\section*{6. Interpretation}

\begin{tabular}{ll}
\hline
$\mathcal{R}$ & Meaning \\
\hline
$1.0 - 1.4$ & Fast bar \\
$> 1.4$ & Slow bar \\
\hline
\end{tabular}

\section*{Conclusion}

The $m=2$ Fourier mode ($A_2$) captures the dominant bar structure in galaxies. 
Its amplitude provides a direct measure of bar strength, while its radial behavior 
allows estimation of the bar length and its dynamical properties.

\begin{thebibliography}{}

    \bibitem{Athanassoula2003}
            Athanassoula, E. 2003, MNRAS, 341, 1179

\bibitem[Laurikainen et al.(2004)]{Laurikainen2004}
Laurikainen, E., Salo, H., \& Buta, R. 2004, ApJ, 607, 103

\bibitem[D{\'\i}az-Garc{\'\i}a et al.(2016)]{DiazGarcia2016}
D{\'\i}az-Garc{\'\i}a, S., Salo, H., Laurikainen, E., \& Herrera-Endoqui, M. 2016, A\&A, 587, A160

\bibitem[Garcia-G{\'o}mez et al.(2017)]{GarciaGomez2017}
Garcia-G{\'o}mez, C., Athanassoula, E., Barber{\`a}, C., \& Bosma, A. 2017, A\&A, 601, A132

\end{thebibliography}

\end{document}
